\section{Tag 5 — Eine neue Mathematik: endliche Körper}

\subsection*{Bleistiftübung 1 — Inverse modulo 10 und 11}

(a) Modulo 10:

\begin{center}
\begin{tabular}{cll}
  \toprule
  $a$ & Inverse & Rechnung \\
  \midrule
  1 & 1   & $1 \cdot 1 = 1$ \\
  2 & --- & keine: $2 \cdot k$ ist immer gerade, kann nie $\equiv 1 \pmod{10}$ sein \\
  3 & 7   & $3 \cdot 7 = 21 \equiv 1$ \\
  4 & --- & gleich wie 2 \\
  5 & --- & $5 \cdot k$ endet auf 0 oder 5 \\
  6 & --- & gleich wie 2 \\
  7 & 3   & $7 \cdot 3 = 21 \equiv 1$ \\
  8 & --- & gleich wie 2 \\
  9 & 9   & $9 \cdot 9 = 81 \equiv 1$ \\
  \bottomrule
\end{tabular}
\end{center}

Inverse haben: 1, 3, 7, 9 — also genau die zu 10 teilerfremden
Zahlen.

(b) Modulo 11:

\begin{center}
\begin{tabular}{cc}
  \toprule
  $a$ & Inverse \\
  \midrule
  1  & 1  \\
  2  & 6  \\
  3  & 4  \\
  4  & 3  \\
  5  & 9  \\
  6  & 2  \\
  7  & 8  \\
  8  & 7  \\
  9  & 5  \\
  10 & 10 \\
  \bottomrule
\end{tabular}
\end{center}

Alle haben Inverse.

(c) 11 ist eine Primzahl, 10 nicht.

(d) \textbf{Faustregel:} eine Zahl $a$ hat genau dann ein Inverses
modulo $n$, wenn $a$ und $n$ teilerfremd sind (d.\,h.\
$\gcd(a, n) = 1$). Folgt daraus: ist $n$ prim, dann sind alle Zahlen
$1, \dots, n-1$ zu $n$ teilerfremd → alle haben Inverse.

\subsection*{Bleistiftübung 2 — GF(2)}

(a) Additionstafel: $0+0=0$, $0+1=1$, $1+0=1$, $1+1=0$.

(b) Multiplikationstafel: $0 \cdot 0=0$, $0 \cdot 1=0$, $1 \cdot 0=0$,
$1 \cdot 1=1$.

(c) Addition entspricht \textbf{XOR}, Multiplikation entspricht
\textbf{AND}.

(d) $1 - 1 = 1 + 1 = 0$ (denn die Subtraktion ist die Addition mit
der Gegenzahl, und in GF(2) ist die Gegenzahl von 1 die 1 selbst).
Allgemeines Prinzip: in GF($2^n$) ist jede Zahl ihre eigene
Gegenzahl.

\subsection*{Bleistiftübung 3 — Modulo 256}

(a) $2 \cdot 128 \bmod 256 = 256 \bmod 256 = 0$. Also: 2 hat kein
Inverses (denn wenn wir mit 2 multiplizieren und etwas erhalten, das
durch 256 teilbar ist, nämlich 0, kann das Ergebnis nicht 1 sein).

(b) Alle \textbf{geraden} Zahlen aus $1, \dots, 255$ haben kein
Inverses, weil sie mit 256 den Faktor 2 gemeinsam haben.

(c) Es gibt 128 ungerade Zahlen ($1, 3, 5, \dots, 255$) — nur diese
sind koprim zu 256.

\subsection*{Bleistiftübung 4 — Addition in GF($2^3$)}

\begin{center}
\begin{tabular}{lll}
  \toprule
  Aufgabe & 3-Bit (XOR) & Polynom \\
  \midrule
  $(x+1) + (x^2+x)$            & $011 \oplus 110 = 101$ & $x^2+1$ \\
  $(x^2+1) + (x^2+x+1)$        & $101 \oplus 111 = 010$ & $x$ \\
  $x^2 + x^2$                  & $100 \oplus 100 = 000$ & $0$ \\
  $(x^2+x+1) + (x^2+x+1)$      & $111 \oplus 111 = 000$ & $0$ \\
  \bottomrule
\end{tabular}
\end{center}

Beobachtung: $a + a = 0$ für jedes Element. Jedes Element ist seine
eigene Gegenzahl.

\subsection*{Bleistiftübung 5 — Multiplikation in GF($2^3$)}

(a) $x \cdot x^2 = x^3$. Reduziere mit $x^3 = x + 1$: Ergebnis =
$\mathbf{x + 1}$ ($= 011$).

(b) $x^2 \cdot x^2 = x^4 = x^2 + x$ ($= 110$).

(c) $(x+1) \cdot (x^2+x)$:
\begin{align*}
  &= x \cdot (x^2+x) + 1 \cdot (x^2+x) \\
  &= x^3+x^2 + x^2+x \\
  &= x^3 + (x^2 + x^2) + x \\
  &= x^3 + 0 + x \\
  &= x^3 + x
\end{align*}
Reduziere $x^3 \to x+1$:
\[
  (x+1) + x = 2x + 1 = 0 \cdot x + 1 = 1
\]
Ergebnis: $\mathbf{1}$. (Das heißt zugleich: $x+1$ und $x^2+x$ sind
zueinander invers.)

(d) $(x^2+x+1) \cdot (x^2+1)$:
\begin{align*}
  &= x^2 \cdot (x^2+1) + x \cdot (x^2+1) + 1 \cdot (x^2+1) \\
  &= x^4+x^2 + x^3+x + x^2+1 \\
  &= x^4 + x^3 + (x^2+x^2) + x + 1 \\
  &= x^4 + x^3 + 0 + x + 1
\end{align*}
Reduziere $x^4 \to x^2+x$ und $x^3 \to x+1$:
\begin{align*}
  &= (x^2+x) + (x+1) + x + 1 \\
  &= x^2 + (x+x+x) + (1+1) \\
  &= x^2 + x + 0 \\
  &= x^2 + x
\end{align*}
Ergebnis: $\mathbf{x^2 + x}$ ($= 110$).

\subsection*{Bleistiftübung 6 — Multiplikationstafel GF($2^3$)}

\begin{center}\small
\begin{tabular}{l|lllllll}
  \toprule
  $\cdot$       & $1$    & $x$    & $x{+}1$ & $x^2$    & $x^2{+}1$ & $x^2{+}x$ & $x^2{+}x{+}1$ \\
  \midrule
  $1$           & $1$    & $x$    & $x{+}1$ & $x^2$    & $x^2{+}1$ & $x^2{+}x$ & $x^2{+}x{+}1$ \\
  $x$           & $x$    & $x^2$  & $x^2{+}x$ & $x{+}1$ & $1$       & $x^2{+}x{+}1$ & $x^2{+}1$ \\
  $x{+}1$       & $x{+}1$ & $x^2{+}x$ & $x^2{+}1$ & $x^2{+}x{+}1$ & $x^2$ & $1$ & $x$ \\
  $x^2$         & $x^2$  & $x{+}1$ & $x^2{+}x{+}1$ & $x^2{+}x$ & $x$    & $x^2{+}1$ & $1$ \\
  $x^2{+}1$     & $x^2{+}1$ & $1$ & $x^2$ & $x$ & $x^2{+}x{+}1$ & $x{+}1$ & $x^2{+}x$ \\
  $x^2{+}x$     & $x^2{+}x$ & $x^2{+}x{+}1$ & $1$ & $x^2{+}1$ & $x{+}1$ & $x$ & $x^2$ \\
  $x^2{+}x{+}1$ & $x^2{+}x{+}1$ & $x^2{+}1$ & $x$ & $1$ & $x^2{+}x$ & $x^2$ & $x{+}1$ \\
  \bottomrule
\end{tabular}
\end{center}

(Einzelne Felder zur Selbstkontrolle: $x \cdot x^2 = x^3 \to x+1$
\checkmark, $(x+1)(x^2+x) = 1$ \checkmark, $(x^2+1) \cdot x = x^3 + x
= (x+1) + x = 1$ \checkmark.)

\subsection*{Bleistiftübung 7 — Inverse aus der Tafel}

\begin{center}
\begin{tabular}{ll}
  \toprule
  Element & Inverses \\
  \midrule
  $1$       & $1$       \\
  $x$       & $x^2+1$   \\
  $x+1$     & $x^2+x$   \\
  $x^2$     & $x^2+x+1$ \\
  $x^2+1$   & $x$       \\
  $x^2+x$   & $x+1$     \\
  $x^2+x+1$ & $x^2$     \\
  \bottomrule
\end{tabular}
\end{center}

Probe: $x \cdot (x^2+1) = x^3 + x = (x+1) + x = 1$ \checkmark.

\subsection*{Aufgaben 1.1, 2.1, 2.2 (Python)}

Erwartete Übereinstimmung mit der Hand-Tafel. Der häufigste
Fehlerquell ist die Reduktionsschleife in \texttt{\_\_mul\_\_} —
wenn die Bit-Indizes nicht stimmen, kommt Müll raus. Beim Debuggen
lohnt es sich, ein konkretes Beispiel von Hand und in Python parallel
laufen zu lassen.

\subsection*{Antworten zu den drei Reflexionsfragen}

\begin{enumerate}
\item \textbf{GF($2^8$) statt GF($2^3$):} vier Änderungen:
  \begin{itemize}
  \item \texttt{MOD} wird zu einem Polynom vom Grad 8, z.\,B.\
    \texttt{0b100011101} (das ist $x^8+x^4+x^3+x^2+1$, das
    standardmäßig für AES und Reed-Solomon im Datamatrix verwendete
    irreduzible Polynom).
  \item Wertebereich-Check $0 \dots 7$ wird zu $0 \dots 255$.
  \item Die Schleife \texttt{for i in range(3)} wird zu \texttt{for
    i in range(8)}.
  \item Die Reduktionsschleife \texttt{for i in range(5, 2, -1)}
    wird zu \texttt{for i in range(14, 7, -1)} (denn das Produkt
    zweier Grad-7-Polynome hat höchstens Grad 14).
  \end{itemize}

\item \textbf{Schnellere Inverse:} eine offensichtliche Idee ist eine
  \textbf{Tabelle} vorzuberechnen (geht bei nur 256 Elementen sehr
  gut — brauche nur 256 Bytes Speicher). Eine andere: den
  \textbf{Satz von Fermat} nutzen, der besagt $a^{2^n - 1} = 1$ in
  GF($2^n$), also $a^{2^n - 2} = a^{-1}$. Für GF($2^8$): $a^{254} =
  a^{-1}$. Das berechnet man mit „schneller Potenzierung“ in nur 8
  Multiplikationen statt 255 Versuchen. Der Erweiterte Euklidische
  Algorithmus ist die mathematisch eleganteste Lösung.

\item Diese Frage ist eine reine Vorschau. Die Idee „Polynome mit
  Koeffizienten aus einem Körper“ ist tatsächlich der Schlüssel zu
  Tag 6: ein RS-Codewort \emph{ist} ein Polynom mit Koeffizienten
  in GF($2^8$).
\end{enumerate}
